\section{Evaluation on Public Benchmarks}
\label{sec:experiemnt}

% 1. Baswed on what insights, we do the following...

% 2. Results (as expected..)

% \zixuan{some may move to introduction}

From the controlled experiments and analysis in \cref{sec:analysis}, we find that MAS
is most effective at the edge of sub-agent competence, instruction-tuned LLMs are better orchestrator initializations than RLMs, and MAS provide consistent gains in {parallel and} adversarial settings. To further validate the effectiveness and generality of \ourframework beyond
synthetic settings, we evaluate it on a range of public  benchmarks spanning diverse domains and task characteristics. 
% \zixuan{what practical guidence you get?}. 
% To futher demonstrate the effectivemnes of \ourframework, we conduct experiments on public benchmarks with various domains.

\textbf{Dataset.} We consider \textbf{5} popular benchmarks, covering math (\textbf{AIME24}, \textbf{AIME25}~\citep{aime2024}), multi-hop QA
% \austin{Is GPQA multi-hop QA? I'd consider it more reasoning QA?} 
(\textbf{GPQA}~\cite{rein2023gpqagraduatelevelgoogleproofqa}, \textbf{HotpotQA}~\citep{DBLP:conf/emnlp/Yang0ZBCSM18}), and multi-step search-based QA (\textbf{BrowseComp+}~\cite{chen2025BrowseCompPlus}). For AIME and GPQA, we use \textbf{DeepScaleR}~\cite{deepscaler2025} as training data. Notably, DeepScaleR mainly focuses on math reasoning and can largely be considered out-of-domain (OOD) for GPQA; thus, GPQA results indicate OOD generalization performance. For HotpotQA, we use the provided training split, and for BrowseComp+, {we use 80\% of the data for training our orchestrator.}
% \shafiq{...}

\textbf{Sub-agents.} For fair comparison, \ourframework uses a fixed set of candidate sub-agents with the same LLM backbone across all sub-agents, varying only their tools and prompting workflows. 
% \ryan{subagent pool is fixed, maybe worth saying explicitly}
We include \textbf{4} widely used workflows {to form sub-agents}: \textbf{CoT} and CoT self-consistency \textbf{(SC)}, \textbf{Debate}, and \textbf{Self-refine}. We also include one \textbf{Search} sub-agent, where the agent automatically decides when and what to search in a multi-turn manner, following OpenDeepResearch \cite{langchain_open_deep_research}. We employ the best setting identified in Section~\ref{sec:analysis}: \qwen as the orchestrator and \ossonetwentylow as the sub-agent backbone.


\textbf{Configurations.} Our analysis shows that MAS should be deployed selectively based on sub-task
structure and underlying LLM capacity, rather than as a default
replacement for a SAS. The proposed \textbf{\masness} notion
allows such task- and model-dependent choices to be explicitly encoded in
\ourframework. Given the largely sequential nature of mathematical and reasoning tasks in
AIME and GPQA, we use \textbf{low \masness}, allowing at most one
sub-agent from \texttt{CoTAgent}, \texttt{SCAgent}, \texttt{DebateAgent}, and \texttt{ReflexionAgent}. In contrast, HotpotQA and BrowseComp+ involve more complex sub-task
structures (e.g., parallel search) and require capabilities that often lie at
the edge of SAS competence (e.g., multi-turn information retrieval). For these
benchmarks, we adopt \textbf{high \masness} and additionally include the \texttt{DeepResearchAgent}.
% , with candidate sub-agents that additionally include a \texttt{DeepResearchAgent}. 
% \update{Specifically, for AIME24, AIME25, and GPQA, we use a low DoM setting with candidate sub-agents including CoT, SC, Debate, and Reflexion. For HotpotQA and BrowseComp+, we adopt a high DoM setting and additionally include a Search sub-agent.}
% \cref{tab:configs} summarizes the
% resulting configurations. 
% \ryan{was \masness for these datasets pre-specified or tuned?}
% Based on the understanding of these benchmarks (see Appendix~\ref{xxx} for more details), 
% we deploy the config below:
% \zixuan{use table}
% \begin{table}[t]
% \setlength{\fboxsep}{1pt}
% \centering
% \resizebox{\columnwidth}{!}{
% \begin{tabular}{ccc}
%\toprule
% \textbf{Benchmark} & \textbf{\masness} & \textbf{Candidate Sub-agents} \\
% \midrule
% AIME24 & Low & CoT, SC, Debate, Reflexion \\
% AIME25 & Low & CoT, SC, Debate, Reflexion \\
% GPQA & Low & CoT, SC, Debate, Reflexion \\
% HotpotQA & High & CoT, SC, Debate, Reflexion, Search \\
% BrowseComp+ & High & CoT, SC, Debate, Reflexion, Search \\
% \bottomrule
% \end{tabular}
% }
% \caption{\small \masness and candidate sub-agents for considered datasets.
% }
% \label{tab:configs}
% \end{table}

\textbf{Overall results.} \cref{tab:overall_public} shows the results against corresponding SAS and other comparable MAS including state-of-the-art (SoTA) \textbf{inference-time orchestration systems} \textbf{AFlow} \citep{zhang2024aflowautomatingagenticworkflow}, \textbf{MaAS} \citep{zhang2025multiagentarchitecturesearchagentic}, \textbf{MAS-Zero} \citep{Ke2025MASZero}, and SoTA public \textbf{training-time orchestration} systems \textbf{MAS-GPT} \citep{ye2025masgpttrainingllmsbuild}, \textbf{ToolOrchestra} \citep{su2025toolorchestraelevatingintelligenceefficient} {(sequential orchestration)}. To the best of our knowledge, these are the only training-time orchestration systems that release trained orchestrators. For a controlled comparison, inference-time baselines use the same orchestrator LLM as \ourframework, while training-time baselines use their officially released trained orchestrators. All systems share \textit{the same candidate sub-agents and sub-agent LLMs} as in \ourframework, {which may differ from their original training environments} (see more details in Sections~\ref{ap:discuss_baselines} and~\ref{ap:addition_observation})

% \shafiq{Give a brief description of those baselines}. 

% \zixuan{TODO: Note that AIME25 is generalized retutls}

% \begin{table}[h]
% % \setlength{\fboxsep}{1pt}
% \centering
% \resizebox{\columnwidth}{!}{
% \begin{tabular}{l cccc c}
% \toprule
% \multirow{2}{*}{\textbf{Method}}
% & \multicolumn{4}{c}{\textbf{IID Tasks}} 
% & \multicolumn{1}{c}{\textbf{OOD Task}} \\
% \cmidrule(lr){2-5}\cmidrule(lr){6-6}
% & \textbf{AIME24} & \textbf{AIME25} & \textbf{HotpotQA} & \textbf{BrowseComp+} & \textbf{GPQA} \\
% \midrule
% CoTAgent       & 50.00 & 45.00 & 33.56 & 1.12 & 60.54 \\
% SCAgent        & 57.50 & 51.67 & 35.50 & 0.75 & 62.88 \\
% DebateAgent    & 62.08 & 57.50 & 36.88 & 0.81 & 64.14 \\
% ReflexionAgent & 60.83 & 50.42 & 36.63 & 1.00 & 62.37 \\
% SearchAgent    & ---   & ---   & 46.44 & 8.56 & ---   \\
% \midrule
% \multicolumn{6}{l}{\textbf{Standalone Agents}} \\
% \midrule
% MaAS    & 32.50  &    &  &  &    \\
% MAS-Zero         & \multicolumn{5}{c}{
% No valid MAS generated with 7B orchestrator, aligned with \citet{Ke2025MASZero}
% } \\
% \midrule
% ToolOrchestra  &  22.92  &    &  &  &    \\
% \midrule
% \textbf{\ourframework} & \textbf{66.25} & \textbf{61.25} & \textbf{49.00} & \textbf{11.00} & \textbf{65.21} \\
% \bottomrule
% \end{tabular}
% }
% \caption{\small \avgat{8} of the considered benchmarks. ``---'' indicates non-applicable. 
% \zixuan{[Working 01/08] lacking baselines?, Training-based baselines include different sub-agents (HotpotQA, BrowseComp+) or datasets. We may only add baselines for AIME and \masness=low}
% % \zixuan{TODO: < 1\% can be removed}
% }
% \label{tab:overall_public}
% \end{table}


\begin{table}[h]
\centering
\setlength{\tabcolsep}{1.5pt}
\resizebox{\columnwidth}{!}{
\begin{tabular}{l cccc c}
\toprule
\multirow{2}{*}{\textbf{Method}}
& \multicolumn{4}{c}{\textbf{IID Tasks}}
& \multicolumn{1}{c}{\textbf{OOD Task}} \\
\cmidrule(lr){2-5}\cmidrule(lr){6-6}
& \textbf{AIME24} & \textbf{AIME25} & \textbf{HotpotQA} & \textbf{BrowseComp+} & \textbf{GPQA} \\
\midrule

\multicolumn{6}{c}{\textbf{Standalone Agents}} \\
\midrule
CoTAgent       & 50.00 & 45.00 & 33.56 & 1.12 & 60.54 \\
SCAgent        & 57.50 & 51.67 & 35.50 & 0.75 & 62.88 \\
DebateAgent    & 62.08 & 57.50 & 36.88 & 0.81 & 64.14 \\
ReflexionAgent & 60.83 & 50.42 & 36.63 & 1.00 & 62.37 \\
DeepResearchAgent    & ---   & ---   & 46.44 & 8.56 & ---   \\
\addlinespace[2pt]

\midrule
\multicolumn{6}{c}{\textbf{SoTA Inference-time Orchestration}} \\
\midrule
AFlow           & 62.50 & 53.33 & --- & --- & 65.43 \\
MaAS           & 32.50 & 40.83 & --- & --- & 40.78 \\
MAS-Zero       & \multicolumn{5}{c}{No valid MAS generated with 7B orchestrator} \\
\addlinespace[2pt]

\midrule
\multicolumn{6}{c}{\textbf{SoTA Public Training-time Orchestration}} \\
\midrule
MAS-GPT   & 58.75 & 43.33 & ---   & ---   & 63.51 \\
ToolOrchestra   & 23.33 & 11.25 & 37.44 &  1.38  & 29.80 \\
\addlinespace[2pt]
\midrule
\multicolumn{6}{c}{\textbf{SoTA LLM as Orchestrator}} \\
\midrule
GPT-5 & 55.00 &	47.72 &	25.87 &  0.50 & 59.01\\
Claude-Sonnet-4.5 & 45.56 &	35.00 & 38.00 &  0.50 & 21.72\\

\addlinespace[2pt]
\midrule
\multicolumn{6}{c}{\textbf{Ours}} \\
\midrule
\textbf{\ourframework} & \textbf{66.25} & \textbf{61.25} & \textbf{49.00} & \textbf{11.00} & \textbf{65.21} \\
\bottomrule
\end{tabular}
}
\caption{\small \avgat{8} accuracy of the considered benchmarks. ``---'' indicates non-applicable.
% \zixuan{4 baselines added. MAS-GPT has SC fallback, so it won't be too bad} 
% \shafiq{reduce the gaps between colomns}\zixuan{reduced}
% \zixuan{[Working 01/08] Lacking baselines? Training-based baselines include different sub-agents (HotpotQA, BrowseComp+) or datasets. We may only add baselines for AIME and \masness = low.}
}
\label{tab:overall_public}
\vspace{-\baselineskip}
\end{table}

% \begin{observation}
% \textbf{\ourframework consistently outperforms the underlying sub-agents {as well as SoTA inference-time and training-time orchestration systems} across all evaluated benchmarks, and it demonstrates strong OOD generalization on GPQA.} 
\textbf{\ourframework consistently outperforms strong baselines across all evaluated benchmarks and demonstrates robust OOD generalization.}
% \shafiq{update based on the newly added baselines}
% We see improvement and align with Sec. 4 observations
% \end{observation}
These results provide concrete evidence of the effectiveness of
\ourframework, 
% Across diverse benchmarks, the framework is able to leverage the
% capabilities of its sub-agents more effectively than any single sub-agent in isolation, 
indicating that performance gains arise from orchestration rather
than from sub-agent strength alone. 
The strong OOD generalization further indicates that the orchestration strategy does not overfit to training distributions, but instead transfers to unseen problems.
% Moreover, the strong OOD generalization shows that the learned orchestration strategy does not overfit, but instead transfers to unseen problem distributions. 
Together, these results confirm that the
principles identified in \cref{sec:analysis} generalize beyond the controlled
setting and translate into consistent improvements on real-world benchmarks.

% also indicates that the findings in \cref{sec:analysis} is generalize beyond the \ourbenchmark. 

% MAS is helpful when parallel, and challenging sub-tasks are invovled. 

% \zixuan{TODO: explian what above ob means and how this show our effectiveness}


\textbf{Further analysis.} We further examine the structure of the MAS designs proposed by
\ourframework. While a comprehensive analysis is provided in the \cref{ap:benchmark_generated_mas}, we
highlight several key findings here that shed light on how the learned
orchestration adapts to different tasks. 


\textbf{Under low \masness, \ourframework learns effective single-agent delegation.}
As shown in \cref{fig:7b_120b_aime24_low}, \ourframework learns to delegate the task entirely to
a single sub-agent (100\% delegation after 20 steps) and dynamically selects
strong sub-agents, primarily \texttt{ReflexionAgent} and \texttt{DebateAgent},
which are the best-performing SAS baselines in
\cref{tab:overall_public}.
This suggests that \ourframework can effectively delegate the task to the best sub-agent under low DoM, which results in its superior performance.
% compared to baselines.


\textbf{Under high \masness, \ourframework learns to exploit parallelism.}
As shown in \cref{fig:7b_120b_browse_comp_plus_high}, \ourframework
learns to
invoke \texttt{DeepResearchAgent} to perform multiple parallel searches,
typically using 3 to 4 per
question.
This results in a characteristic MAS pattern on BrowseComp+, where relevant information is first collected through parallel search and then combined by an aggregation agent. Such a global orchestration strategy underlies the improved performance.
% {This reflects a typical MAS pattern produced by \ourframework on BrowseComp+: the system first retrieves potentially relevant information through multiple parallel searches and then introduces an aggregation agent to combine the retrieved results. This behavior is effective and highlights the global view adopted by \ourframework during MAS design: collecting comprehensive evidence before aggregation improves final answer correctness.}

% This indicate the typical pattern of generated MAS for BrowseComp+ first retrieves all potentially required information through multiple search agents and then introduces an organization or aggregation agent to combine the retrieved results. This behavior is effcive and clearly show the global view when \ourframework design the MAS: collecting comprehensive evidence before aggregation can improve final-answer correctness.

\textbf{\ourframework achieves performance–cost Pareto frontier.}
As shown in \cref{fig:parato_front}, \ourframework lies on the Pareto frontier among the considered inference-time and training-time orchestration systems, achieving higher accuracy at lower or comparable cost. This efficiency follows directly from the behaviors observed earlier: \ourframework dynamically adapts to each task, generating MAS designs that match the underlying sub-task structure and delegating execution to the most effective agent configurations.

% \textbf{\ourframework achieves perforamnce-cost parato-frnt. }As shwon in \cref{fig:parato_front}, \ourframework lies in the pareofron acorss the consdierd inference-time and training-time orchestation systems, achieving higher accuracy at lower or comparable cost. This is achieveable as \ourframework dynamically adapts to the given task by
% proposing MAS designs that align with the underlying sub-task structure and by
% delegating execution to the most effective agent configurations, as we have seen in the previoous observatiopn.

% \cref{tab:cost_breakdown}

% \zixuan{01/21 efficiency advantage? a subsec?}
% \zixuan{below pargraph is not neccessary}

% {Taken together, the observations under low and high \masness indicate that \ourframework dynamically adapts to the given task by
% proposing MAS designs that align with the underlying sub-task structure and by
% delegating execution to the most effective agent configurations.}



% \begin{discussion}
% \textbf{Training dynamic and Proposed MAS.}
% \zixuan{MISSING 6: statistic, dynamix and observation}

% \end{discussion}

% \begin{takeaway}
% \end{takeaway}
